\chapter{Literature Review}

\setstretch{2}
{\fontspec{Times New Roman}\fontsize{12pt}{14pt}\selectfont
	
	\section{Scope and Positioning of This Review}
	This chapter reviews prior work from a \textbf{project-construction perspective} for RSHub-agent, rather than from a single-feature perspective. The goal is to identify which technical lines are required to build an end-to-end remote sensing agent system, including: (1) multi-modal remote sensing task characteristics, (2) LLM-agent architectural choices, and (3) key enabling mechanisms such as HITL, RAG, and layered skill loading.

	\section{Remote Sensing Workflow Requirements}
	
	\subsection{Multi-Modal Data and Processing Constraints}
	Remote sensing applications commonly combine optical, SAR, and polarimetric signals, which introduces heterogeneity in spatial resolution, noise characteristics, and feature representation [30][32][33]. For practical systems, this means an agent must support robust parameter handling and cross-source workflow coordination instead of relying on single-step prompt response.
	
	\subsection{Application Characteristics}
	Recent systems such as PolSAM [30], SkySense [33], MineAgent [32], Tree-GPT [37], and Earth-Agent [50] demonstrate that real tasks are often multi-stage (task planning, parameter completion, execution, and result interpretation) and require stable context tracking across turns.

	\section{LLM-driven Agent Architectures}
	
	\subsection{Multi-Agent Architectures}
	Multi-agent frameworks decompose complex tasks into specialized sub-agents and often improve parallelism [1][35][36]. GeoLLM-Squad [1] and EarthGPT-X [35] show strong capability in heterogeneous sensing scenarios. However, communication and synchronization overhead can increase engineering complexity, especially when workflow control and state consistency are critical.
	
	\subsection{Single-Agent Architectures}
	Single-agent systems integrate routing, reasoning, and tool invocation in one orchestrator. RS-Agent [3] and EarthMarker [26] show that modular single-agent designs can maintain good performance with lower coordination overhead. For project-oriented deployment, this paradigm is easier to integrate with existing APIs, session state, and billing controls.
	
	\subsection{Comparative Analysis}
	Overall, multi-agent designs are strong in parallel decomposition, while single-agent designs are typically stronger in operational simplicity and maintainability [37][50]. Table~\ref{tab:agent-architecture} summarizes representative trade-offs.
	
	\begin{table}[h!]
		\centering
		\footnotesize
		\caption{Comparison of Multi-Agent and Single-Agent Architectures}
		\label{tab:agent-architecture}
		\begin{tabular}{lcccc}
			\hline
			System & Architecture & Parallelism & Memory & Notes \\
			\hline
			GeoLLM-Squad [1] & Multi-agent & High & Medium & Domain-specific agents \\
			RS-Agent [3] & Single-agent & Medium & Low & Modular single-agent \\
			SkyMoE [27] & Multi-agent & High & High & Mixture-of-experts, multi-modal \\
			EarthMarker [26] & Single-agent & Medium & Medium & Region-/point-level prompts \\
			EarthGPT-X [35] & Multi-agent & High & High & Flexible multi-source analysis \\
			Tree-GPT [37] & Single-agent & Medium & Medium & Forest image understanding \\
			MineAgent [32] & Single-agent & Medium & Medium & Mineral exploration analysis \\
			Earth-Agent [50] & Single-agent & Medium & Medium & Comprehensive Earth observation \\
			\hline
		\end{tabular}
	\end{table}
	
	\section{Key Enabling Mechanisms for Project-level Agents}
	
	\subsection{Human-in-the-Loop (HITL)}
	HITL injects explicit human checkpoints into automated pipelines to improve safety and controllability [14][47]. In remote sensing task submission scenarios, the Plan $\rightarrow$ Check $\rightarrow$ Confirm $\rightarrow$ Execute pattern is widely used to prevent unintended execution and error propagation [36][48].
	
	\begin{table}[h!]
		\centering
		\footnotesize
		\caption{HITL Framework Key Steps}
		\label{tab:hitl-steps}
		\begin{tabular}{l p{8cm} p{4cm}}
			\hline
			Step & Description & Benefit \\
			\hline
			Plan & Parse user request into structured task spec & Ensure task completeness \\
			Check & Verify parameters and constraints & Prevent execution errors \\
			Confirm & Obtain human approval & Increase reliability \\
			Execute & Run task after confirmation & Maintain correctness and reproducibility \\
			\hline
		\end{tabular}
	\end{table}
	
	\subsection{Retrieval-Augmented Generation (RAG)}
	RAG improves factual grounding by retrieving task-relevant knowledge during generation [12][13][31][49]. In long multi-turn interactions, RAG helps reduce context drift and improves instruction relevance compared with static monolithic prompts [29][36][37].
	
	\subsection{Skill Architecture and Layered Loading}
	Layered skill systems separate lightweight metadata (catalog/index layer) from detailed full documents (execution layer), so that only relevant instructions are injected per turn [5][7][26]. This design improves context efficiency and supports scalable expansion when new domain tools are introduced.
	
	\begin{figure}[h!]
		\centering
		\begin{tikzpicture}[
			node distance=1.5cm,
			box/.style={rectangle, draw, rounded corners, minimum width=5cm, minimum height=0.8cm, align=center, fill=#1!20},
			arrow/.style={-{Stealth[scale=1.0]}, thick}
			]
			
			\node[box=blue] (layer1) {Layer 1\\Skill Catalog\\(Lightweight)};
			\node[box=green, below=of layer1] (layer2) {Layer 2\\Full Skill Documents\\(On-demand)};
			\node[box=orange, below=of layer2] (rag) {RAG-enhanced\\Selection};
			\node[box=red, below=of rag] (task) {Task Execution};
			
			\draw[arrow] (layer1) -- (layer2);
			\draw[arrow] (layer2) -- (rag);
			\draw[arrow] (rag) -- (task);
			
			\node[left=0.5cm of layer1] {Lightweight retrieval, minimal token usage};
			\node[left=0.5cm of layer2] {Detailed skill info, precise injection};
			\node[left=0.5cm of rag] {Dynamic retrieval of relevant knowledge};
			\node[left=0.5cm of task] {Execution of the selected task};
			
		\end{tikzpicture}
		\caption{Conceptual diagram of RAG integration with Layered Skill architecture.}
		\label{fig:rag-layered-tikz-vertical}
	\end{figure}
	
	\section{Summary and Research Gaps}
	
	The literature indicates that building a practical remote sensing agent requires coordinated design across architecture, safety control, retrieval, and skill organization. Multi-agent systems offer stronger parallelism, while modular single-agent systems are often easier to deploy and maintain in project settings [3][26][50].
	
	For RSHub-agent, three project-level gaps remain important: (1) balancing HITL safety with interaction efficiency, (2) reducing long-context overhead while preserving response quality, and (3) improving engineering maintainability for continuous updates and synchronization. These gaps motivate the system design, software specification, and validation strategy presented in subsequent chapters.
	
}